\section{The characteristic equation}
We use the book's convention $p(\lambda)=\det(A-\lambda I)$.
Some software instead reports $\det(\lambda I-A)$; they differ
by $(-1)^n$ and have the same zeros. An eigenvalue is a zero of
$p$, and its algebraic multiplicity is its multiplicity as a polynomial
root. That multiplicity need not equal the dimension of its eigenspace.
For a triangular matrix, the characteristic determinant is the product
of the shifted diagonal entries. For a general matrix, use the full
determinant rather than multiplying the original diagonal entries.

\practice{5.2.P1}{A characteristic equation with no real roots}
For the given matrix,
\[
 p(\lambda)=\det\begin{bmatrix}1-\lambda&-4\\4&2-\lambda\end{bmatrix}
 =(1-\lambda)(2-\lambda)+16=\lambda^2-3\lambda+18.
\]
Its discriminant is $9-72=-63<0$, so the characteristic equation
has no real root and $A$ has no real eigenvalue. If complex scalars
are allowed, its eigenvalues are $(3\pm i\sqrt{63})/2$.

\exercise{5.2.1}{Compute a two-by-two characteristic polynomial}
For $A=\begin{bmatrix}2 & 7\\7 & 2\end{bmatrix}$, subtract $\lambda$ only from the diagonal entries. The determinant formula $ad-bc$ gives
\[ p(\lambda)=(2 - \lambda)(2 - \lambda)-(7)(7)=\lambda^{2} - 4 \lambda - 45.\]

Factoring gives $p(\lambda)=\left(\lambda - 9\right) \left(\lambda + 5\right)$. Its real roots, hence the eigenvalues, are $-5,\ 9$. Each listed root has multiplicity one.

\exercise{5.2.2}{Compute a two-by-two characteristic polynomial}
For $A=\begin{bmatrix}-4 & -1\\6 & 1\end{bmatrix}$, subtract $\lambda$ only from the diagonal entries. The determinant formula $ad-bc$ gives
\[ p(\lambda)=(- \lambda - 4)(1 - \lambda)-(-1)(6)=\lambda^{2} + 3 \lambda + 2.\]

Factoring gives $p(\lambda)=\left(\lambda + 1\right) \left(\lambda + 2\right)$. Its real roots, hence the eigenvalues, are $-2,\ -1$. Each listed root has multiplicity one.

\exercise{5.2.3}{Compute a two-by-two characteristic polynomial}
For $A=\begin{bmatrix}-4 & 2\\6 & 7\end{bmatrix}$, subtract $\lambda$ only from the diagonal entries. The determinant formula $ad-bc$ gives
\[ p(\lambda)=(- \lambda - 4)(7 - \lambda)-(2)(6)=\lambda^{2} - 3 \lambda - 40.\]

Factoring gives $p(\lambda)=\left(\lambda - 8\right) \left(\lambda + 5\right)$. Its real roots, hence the eigenvalues, are $-5,\ 8$. Each listed root has multiplicity one.

\exercise{5.2.4}{Compute a two-by-two characteristic polynomial}
For $A=\begin{bmatrix}8 & 2\\3 & 3\end{bmatrix}$, subtract $\lambda$ only from the diagonal entries. The determinant formula $ad-bc$ gives
\[ p(\lambda)=(8 - \lambda)(3 - \lambda)-(2)(3)=\lambda^{2} - 11 \lambda + 18.\]

Factoring gives $p(\lambda)=\left(\lambda - 9\right) \left(\lambda - 2\right)$. Its real roots, hence the eigenvalues, are $2,\ 9$. Each listed root has multiplicity one.

\exercise{5.2.5}{Compute a two-by-two characteristic polynomial}
For $A=\begin{bmatrix}8 & 4\\4 & 8\end{bmatrix}$, subtract $\lambda$ only from the diagonal entries. The determinant formula $ad-bc$ gives
\[ p(\lambda)=(8 - \lambda)(8 - \lambda)-(4)(4)=\lambda^{2} - 16 \lambda + 48.\]

Factoring gives $p(\lambda)=\left(\lambda - 12\right) \left(\lambda - 4\right)$. Its real roots, hence the eigenvalues, are $4,\ 12$. Each listed root has multiplicity one.

\exercise{5.2.6}{Compute a two-by-two characteristic polynomial}
For $A=\begin{bmatrix}9 & -2\\2 & 5\end{bmatrix}$, subtract $\lambda$ only from the diagonal entries. The determinant formula $ad-bc$ gives
\[ p(\lambda)=(9 - \lambda)(5 - \lambda)-(-2)(2)=\lambda^{2} - 14 \lambda + 49.\]

Factoring gives $p(\lambda)=\left(\lambda - 7\right)^{2}$. Its real roots, hence the eigenvalues, are $7$. The root $7$ has multiplicity two.

\exercise{5.2.7}{Compute a two-by-two characteristic polynomial}
For $A=\begin{bmatrix}5 & 3\\-4 & 4\end{bmatrix}$, subtract $\lambda$ only from the diagonal entries. The determinant formula $ad-bc$ gives
\[ p(\lambda)=(5 - \lambda)(4 - \lambda)-(3)(-4)=\lambda^{2} - 9 \lambda + 32.\]

The discriminant is $(-9)^2-4(32)=-47<0$. Consequently there are no real eigenvalues. The two complex roots, if desired, are $(9\pm i\sqrt{47})/2$.

\exercise{5.2.8}{Compute a two-by-two characteristic polynomial}
For $A=\begin{bmatrix}-4 & 3\\2 & 1\end{bmatrix}$, subtract $\lambda$ only from the diagonal entries. The determinant formula $ad-bc$ gives
\[ p(\lambda)=(- \lambda - 4)(1 - \lambda)-(3)(2)=\lambda^{2} + 3 \lambda - 10.\]

Factoring gives $p(\lambda)=\left(\lambda - 2\right) \left(\lambda + 5\right)$. Its real roots, hence the eigenvalues, are $-5,\ 2$. Each listed root has multiplicity one.

\exercise{5.2.9}{Expand the shifted determinant using its zeros}
The matrix whose determinant is required is
\[ A-\lambda I=\begin{bmatrix}4 - \lambda & 0 & -1\\0 & 4 - \lambda & -1\\1 & 0 & 2 - \lambda\end{bmatrix}.\]

Expand down column two; only its middle entry is nonzero. Thus
\[ p(\lambda)=(4-\lambda)[(4-\lambda)(2-\lambda)+1].\]

Multiplying the factors gives
\[ \boxed{p(\lambda)=- \lambda^{3} + 10 \lambda^{2} - 33 \lambda + 36}.\]

A useful factored check is $p(\lambda)=- \left(\lambda - 4\right) \left(\lambda - 3\right)^{2}$. Expanding this expression reproduces the same cubic and, at $\lambda=0$, the original determinant.

\exercise{5.2.10}{Expand the shifted determinant using its zeros}
The matrix whose determinant is required is
\[ A-\lambda I=\begin{bmatrix}3 - \lambda & 1 & 1\\0 & 5 - \lambda & 0\\-2 & 0 & 7 - \lambda\end{bmatrix}.\]

Expand across row two; only its middle entry is nonzero. Thus
\[ p(\lambda)=(5-\lambda)[(3-\lambda)(7-\lambda)+2].\]

Multiplying the factors gives
\[ \boxed{p(\lambda)=- \lambda^{3} + 15 \lambda^{2} - 73 \lambda + 115}.\]

A useful factored check is $p(\lambda)=- \left(\lambda - 5\right) \left(\lambda^{2} - 10 \lambda + 23\right)$. Expanding this expression reproduces the same cubic and, at $\lambda=0$, the original determinant.

\exercise{5.2.11}{Expand the shifted determinant using its zeros}
The matrix whose determinant is required is
\[ A-\lambda I=\begin{bmatrix}3 - \lambda & 0 & 0\\2 & 1 - \lambda & 4\\1 & 0 & 4 - \lambda\end{bmatrix}.\]

Expand down column two; only its middle entry is nonzero, and the remaining minor is triangular. Thus
\[ p(\lambda)=(1-\lambda)(3-\lambda)(4-\lambda).\]

Multiplying the factors gives
\[ \boxed{p(\lambda)=- \lambda^{3} + 8 \lambda^{2} - 19 \lambda + 12}.\]

A useful factored check is $p(\lambda)=- \left(\lambda - 4\right) \left(\lambda - 3\right) \left(\lambda - 1\right)$. Expanding this expression reproduces the same cubic and, at $\lambda=0$, the original determinant.

\exercise{5.2.12}{Expand the shifted determinant using its zeros}
The matrix whose determinant is required is
\[ A-\lambda I=\begin{bmatrix}- \lambda - 1 & 0 & 2\\3 & 1 - \lambda & 0\\0 & 1 & 2 - \lambda\end{bmatrix}.\]

Expand across row one. The middle term is zero; the last cofactor contributes $2(3)=6$. Thus
\[ p(\lambda)=(-1-\lambda)(1-\lambda)(2-\lambda)+6.\]

Multiplying the factors gives
\[ \boxed{p(\lambda)=- \lambda^{3} + 2 \lambda^{2} + \lambda + 4}.\]

Factoring is not required here. As checks, the leading term is $-\lambda^3$ under our convention, and substituting $\lambda=0$ gives $\det A=4$.

\exercise{5.2.13}{Expand the shifted determinant using its zeros}
The matrix whose determinant is required is
\[ A-\lambda I=\begin{bmatrix}6 - \lambda & -2 & 0\\-2 & 9 - \lambda & 0\\5 & 8 & 3 - \lambda\end{bmatrix}.\]

Expand down column three; only its last entry is nonzero. Thus
\[ p(\lambda)=(3-\lambda)[(6-\lambda)(9-\lambda)-4].\]

Multiplying the factors gives
\[ \boxed{p(\lambda)=- \lambda^{3} + 18 \lambda^{2} - 95 \lambda + 150}.\]

A useful factored check is $p(\lambda)=- \left(\lambda - 10\right) \left(\lambda - 5\right) \left(\lambda - 3\right)$. Expanding this expression reproduces the same cubic and, at $\lambda=0$, the original determinant.

\exercise{5.2.14}{Expand the shifted determinant using its zeros}
The matrix whose determinant is required is
\[ A-\lambda I=\begin{bmatrix}4 - \lambda & 0 & -1\\-1 & - \lambda & 4\\0 & 2 & 3 - \lambda\end{bmatrix}.\]

Expand across row one. The first minor gives $(-\lambda)(3-\lambda)-8$, and the last cofactor contributes $(-1)(-2)=2$. Thus
\[ p(\lambda)=(4-\lambda)[(-\lambda)(3-\lambda)-8]+2.\]

Multiplying the factors gives
\[ \boxed{p(\lambda)=- \lambda^{3} + 7 \lambda^{2} - 4 \lambda - 30}.\]

A useful factored check is $p(\lambda)=- \left(\lambda - 5\right) \left(\lambda^{2} - 2 \lambda - 6\right)$. Expanding this expression reproduces the same cubic and, at $\lambda=0$, the original determinant.

\exercise{5.2.15}{Read triangular eigenvalues with repetition}
This matrix is upper triangular. Therefore $p(\lambda)$ is the product of its diagonal entries minus $\lambda$:
\[ p(\lambda)=\left(\lambda - 5\right)^{2} \left(\lambda - 3\right) \left(\lambda - 2\right).\]

Listing the diagonal entries, with each occurrence retained, gives the requested eigenvalue list $2,\ 3,\ 5,\ 5$. Repeated entries are repeated eigenvalues in this list, not extra distinct values.

\exercise{5.2.16}{Read triangular eigenvalues with repetition}
This matrix is lower triangular. Therefore $p(\lambda)$ is the product of its diagonal entries minus $\lambda$:
\[ p(\lambda)=\left(\lambda - 6\right) \left(\lambda - 3\right) \left(\lambda - 2\right) \left(\lambda + 5\right).\]

Listing the diagonal entries, with each occurrence retained, gives the requested eigenvalue list $-5,\ 2,\ 3,\ 6$. Repeated entries are repeated eigenvalues in this list, not extra distinct values.

\exercise{5.2.17}{Read triangular eigenvalues with repetition}
This matrix is lower triangular. Therefore $p(\lambda)$ is the product of its diagonal entries minus $\lambda$:
\[ p(\lambda)=- \lambda \left(\lambda - 3\right)^{2} \left(\lambda - 1\right)^{2}.\]

Listing the diagonal entries, with each occurrence retained, gives the requested eigenvalue list $0,\ 1,\ 1,\ 3,\ 3$. Repeated entries are repeated eigenvalues in this list, not extra distinct values.


\exercise{5.2.18}{Force a two-dimensional eigenspace for four}
Subtract $4I$ to obtain
\[
 A-4I=\begin{bmatrix}0&2&3&3\\0&-2&h&3\\0&0&0&14\\0&0&0&-2\end{bmatrix}.
\]
The last row of $(A-4I)x=0$ forces $x_4=0$. The first row then
gives $2x_2+3x_3=0$, so $x_2=-3x_3/2$. Substituting into
the second gives $(h+3)x_3=0$. The coordinate $x_1$ is always free.
For the eigenspace to have a second free coordinate, we need
$h+3=0$. At $h=-3$ its basis can be chosen as
\[
 (1,0,0,0)^T,\qquad(0,-3,2,0)^T.
\]
For any other $h$, $x_3=x_2=x_4=0$, leaving only one free coordinate.
\answer{$\boxed{h=-3}$.}

\exercise{5.2.19}{The constant coefficient is the determinant}
The supplied factorization is
$p(\lambda)=\prod_{j=1}^n(\lambda_j-\lambda)$, with eigenvalues
repeated according to multiplicity. Substitute $\lambda=0$:
\[
 \det A=\det(A-0I)=p(0)=\lambda_1\cdots\lambda_n.
\]
No sign change is needed under this convention; using
$\det(\lambda I-A)$ would introduce the corresponding constant-term sign.
The same evaluation works over complex numbers when nonreal roots are present.

\exercise{5.2.20}{Transpose preserves the whole characteristic polynomial}
Since $I^T=I$ and transposition commutes with subtraction,
$A^T-\lambda I=(A-\lambda I)^T$.
The determinant of a transpose equals the original determinant. Hence
\[
 \det(A^T-\lambda I)=\det((A-\lambda I)^T)=\det(A-\lambda I)
\]
for every $\lambda$. Thus the polynomials are identical, so the
roots agree with their algebraic multiplicities, not merely as distinct sets.

\exercise{5.2.21}{Determinants and the sign of a root}
\begin{enumerate}[label=(\alph*)]
\item \textbf{False.} For $A=\begin{bmatrix}0&1\\1&0\end{bmatrix}$,
$\det A=-1$, but the product of the diagonal entries is $0$.
The diagonal-product rule is valid for triangular matrices.
\item \textbf{False.} A row interchange negates the determinant;
scaling a row by $c$ multiplies it by $c$. Only row replacement
preserves it unchanged in general.
\item \textbf{True.} This is the determinant product law for square matrices.
\item \textbf{False.} A factor $\lambda+5$ vanishes at $\lambda=-5$,
not $5$. For example $A=[-5]$ has that factor up to the harmless
leading sign, but has only eigenvalue $-5$.
\end{enumerate}

\exercise{5.2.22}{Volume, transpose, multiplicity, and row operations}
\begin{enumerate}[label=(\alph*)]
\item \textbf{False as stated.} The geometric volume is $|\det A|$.
The determinant can be negative because it also records orientation.
\item \textbf{False.} $\det A^T=\det A$, with no factor $-1$.
The identity matrix gives $1$, not $-1$.
\item \textbf{True.} This is the definition of algebraic multiplicity.
\item \textbf{False.} Start with
$\begin{bmatrix}1&1\\1&1\end{bmatrix}$, whose eigenvalues are
$2,0$. Replacing row two by row two minus row one gives
$\begin{bmatrix}1&1\\0&0\end{bmatrix}$, whose eigenvalues are
$1,0$. Row replacement preserves the determinant, not every eigenvalue.
\end{enumerate}

\exercise{5.2.23}{Interchanging two factors gives a similar matrix}
Assume $A=QR$ with $Q$ invertible. Then
\[
 Q^{-1}AQ=Q^{-1}(QR)Q=(Q^{-1}Q)RQ=RQ=A_1.
\]
Thus $A_1=Q^{-1}AQ$, precisely the definition of similarity.
Only $Q$ needs to be invertible; $R$ need not be.
This is why interchanging the factors in the QR iteration preserves eigenvalues.

\exercise{5.2.24}{Similarity preserves determinants}
If $B=P^{-1}AP$ with $P$ invertible, the product law gives
\[
 \det B=\det(P^{-1})\det A\det P.
\]
Since $\det(P^{-1})\det P=\det I=1$, the two outer factors
cancel and $\det B=\det A$. This proof also covers singular $A$:
only $P$, not $A$, is required to have an inverse.

\exercise{5.2.25}{Resolve a probability vector into two eigenvectors}
\step{(a) Find a second eigendirection.}
The stated $v_1=(3/7,4/7)^T$ satisfies $Av_1=v_1$.
For $v_2=(1,-1)^T$, multiplication gives
$Av_2=(0.3,-0.3)^T=0.3v_2$. The vectors correspond to distinct
eigenvalues $1$ and $0.3$, so they are independent and form a basis of $\R^2$.
\step{(b) Expand the initial vector.}
\[
 x_0-v_1=(1/2-3/7,1/2-4/7)^T=(1/14,-1/14)^T=\tfrac1{14}v_2.
\]
Thus $x_0=v_1+(1/14)v_2$.
\step{(c) Use the eigenvalue powers.}
Direct multiplication gives $x_1=(0.45,0.55)^T$ and
$x_2=(0.435,0.565)^T$. For every $k\ge0$,
\[
 x_k=A^kx_0=v_1+\frac1{14}(0.3)^kv_2.
\]
Because $|0.3|<1$, the scalar factor tends to zero, so $x_k\to v_1$.
This proves convergence, rather than merely inferring it from two similar iterates.

\exercise{5.2.26}{Derive the two-by-two determinant from elimination}
\step{Case 1: $a\ne0$.}
Replace row two by row two minus $(c/a)$ times row one. This operation
does not change the determinant, and it produces
$\begin{bmatrix}a&b\\0&d-cb/a\end{bmatrix}$.
The triangular formula gives $a(d-cb/a)=ad-bc$.
\step{Case 2: $a=0$.}
If also $c=0$, the first column is zero, so the determinant is zero,
agreeing with $ad-bc$. If $c\ne0$, interchange rows to produce
$\begin{bmatrix}c&d\\0&b\end{bmatrix}$. Its determinant is $cb$,
and one interchange reverses the sign, so the original determinant
is $-cb=ad-bc$. These cases include all possible values of $a,c$.

\exercise{5.2.27}{Three eigendirections determine a steady-state limit}
\step{(a) Check the three given images.}
With the matrix in the exercise, direct multiplication gives
\[
 Av_1=(0.3,0.6,0.1)^T=v_1,
\]
\[
 Av_2=(0.5,-1.5,1)^T=\tfrac12v_2,\qquad
 Av_3=(-0.2,0,0.2)^T=\tfrac15v_3.
\]
All three vectors are nonzero, so these really are eigenvector equations.
\step{(b) Use both a basis and the coordinate sum.}
The eigenvalues $1,1/2,1/5$ are distinct, so the three eigenvectors
are independent and form a basis of $\R^3$. Thus any $x_0$ has
unique coefficients $c_j$ in $x_0=c_1v_1+c_2v_2+c_3v_3$.
For $w=(1,1,1)^T$, the probability hypothesis gives $w^Tx_0=1$.
But $w^Tv_1=1$, $w^Tv_2=0$, and $w^Tv_3=0$. Applying
$w^T$ to the expansion therefore yields $1=c_1$.
\step{(c) Let the contracting components vanish.}
\[
 x_k=A^kx_0=v_1+c_2(1/2)^kv_2+c_3(1/5)^kv_3.
\]
The two powers tend to zero, so $x_k\to v_1$.
The two coefficients may depend on the initial distribution, but the
limit does not, because its coefficient was fixed by total probability one.

\exercise{5.2.28}{Reproducible random examples and the transpose distinction}
To make the experiment reproducible, draw upper-triangular entries
uniformly from the integers $-3,\ldots,4$ with pseudorandom seed
$52028$, setting entries below the diagonal to zero. Successive draws
of sizes four and five give the following valid random integer examples.
The triangular choice makes exact checks of every multiplicity transparent.


\[ A_4=\begin{bmatrix}-1 & 0 & 0 & 2\\0 & 1 & -2 & 1\\0 & 0 & 1 & 4\\0 & 0 & 0 & 2\end{bmatrix}.\]

Exact calculation gives the same polynomial for $A_4$ and $A_4^T$:
\[ p(\lambda)=\left(\lambda - 2\right) \left(\lambda - 1\right)^{2} \left(\lambda + 1\right).\]

Its eigenvalues, with multiplicity, are $-1,\ 1,\ 1,\ 2$. However $e_1$ is an eigenvector of $A_4$ with eigenvalue $-1$, whereas $A_4^Te_1=\begin{bmatrix}-1\\0\\0\\2\end{bmatrix}$ is not a multiple of $e_1$. Thus their eigenvectors are not the same.


\[ A_5=\begin{bmatrix}1 & 3 & 2 & 3 & 0\\0 & -3 & -3 & 4 & -3\\0 & 0 & 4 & -3 & 4\\0 & 0 & 0 & 4 & 3\\0 & 0 & 0 & 0 & 0\end{bmatrix}.\]

Exact calculation gives the same polynomial for $A_5$ and $A_5^T$:
\[ p(\lambda)=- \lambda \left(\lambda - 4\right)^{2} \left(\lambda - 1\right) \left(\lambda + 3\right).\]

Its eigenvalues, with multiplicity, are $-3,\ 0,\ 1,\ 4,\ 4$. However $e_1$ is an eigenvector of $A_5$ with eigenvalue $1$, whereas $A_5^Te_1=\begin{bmatrix}1\\3\\2\\3\\0\end{bmatrix}$ is not a multiple of $e_1$. Thus their eigenvectors are not the same.

The equality of characteristic polynomials is guaranteed in general by Exercise 20; equal eigenspaces are not guaranteed. If a particular matrix is symmetric, its transpose is identical and its eigenspaces do agree, so one must not claim they always differ.


\exercise{5.2.29}{Compare a determinant with pivot and eigenvalue products}
One reproducible choice of a random integer matrix is
\[
 A=\begin{bmatrix}2&3&-3&-2\\5&-3&0&3\\-3&-4&0&2\\1&1&4&1\end{bmatrix}.
\]
It was selected from pseudorandom entries in $-4,\ldots,5$
(seed $52029$) with nonzero pivots and real eigenvalues, so every
requested quantity can be displayed in real decimal form.
\step{(a) Eliminate without row scaling.}
First use $R_2\leftarrow R_2-(5/2)R_1$,
$R_3\leftarrow R_3+(3/2)R_1$, and
$R_4\leftarrow R_4-(1/2)R_1$. Continue replacing rows to clear
entries below the second and third pivots. No swaps or row scalings
are needed. The resulting echelon matrix is
\[
 U=\begin{bmatrix}2&3&-3&-2\\0&-21/2&15/2&8\\0&0&-29/7&-13/21\\0&0&0&74/87\end{bmatrix}.
\]
Row replacement preserves the determinant, so
\[
 \det A=2(-21/2)(-29/7)(74/87)=74.
\]
\step{(b) Compute and multiply the eigenvalues.}
Exact computation gives $p(\lambda)=\lambda^4-40\lambda^2-22\lambda+74$.
Its four real roots, computed with guard digits, are approximately
\[
 -5.8358727568,\quad -1.7435586500,\quad
 1.1271321671,\quad6.4522992397.
\]
Their product using unrounded roots is $74$ to numerical precision,
as also follows exactly from the polynomial's constant coefficient.
\step{(c) Give the requested four-decimal report and comparison.}
The pivots are $2.0000,-10.5000,-4.1429,0.8506$; the eigenvalues
are $-5.8359,-1.7436,1.1271,6.4523$.
The direct matrix determinant is exactly $74$, agreeing with both
products before rounding. Products formed only from the four-place
display may differ slightly because the displayed numbers have been rounded.

\exercise{5.2.30}{A small entry change can destroy two real eigenvalues}
For the stated matrix, expansion of the shifted determinant gives
\[
 p_a(t)=-t^3+4t^2+(379-12a)t+(12a-382)
       =-(t-1)(t^2-3t+12a-382).
\]
Thus $1$ is always an eigenvalue, and the other two are
$(3\pm\sqrt{1537-48a})/2$. Substituting each requested value gives
\[
\begin{array}{c|l|l}
 a&p_a(t)&\text{eigenvalues (approximately)}\\\hline
32&-t^3+4t^2-5t+2&1,1,2\\
31.9&-t^3+4t^2-3.8t+0.8&0.295841,1,2.704159\\
31.8&-t^3+4t^2-2.6t-0.4&-0.127882,1,3.127882\\
32.1&-t^3+4t^2-6.2t+3.2&1,1.5\pm0.974679i\\
32.2&-t^3+4t^2-7.4t+4.4&1,1.5\pm1.466288i
\end{array}
\]
The complex entries in the last two rows are the two nonreal eigenvalues;
the only real eigenvalue in those cases is $1$.
Here are all five graphs on the requested common interval $0\le t\le3$.

\begin{center}\setlength{\unitlength}{1mm}\begin{picture}(145,90)
\put(15,50){\vector(1,0){118}}\put(20,10){\vector(0,1){70}}
\put(135,47){$t$}\put(3,81){$p_a(t)$}

\put(20,49){\line(0,1){2}}\put(20,45){\makebox(0,0){\scriptsize$0$}}

\put(55,49){\line(0,1){2}}\put(55,45){\makebox(0,0){\scriptsize$1$}}

\put(90,49){\line(0,1){2}}\put(90,45){\makebox(0,0){\scriptsize$2$}}

\put(125,49){\line(0,1){2}}\put(125,45){\makebox(0,0){\scriptsize$3$}}

\put(19,18){\line(1,0){2}}\put(14,18){\makebox(0,0){\scriptsize$-8$}}

\put(19,34){\line(1,0){2}}\put(14,34){\makebox(0,0){\scriptsize$-4$}}

\put(19,66){\line(1,0){2}}\put(14,66){\makebox(0,0){\scriptsize$4$}}

{\color{blue}\linethickness{0.8pt}\polyline(20.000,58.000)(21.167,57.351)(22.333,56.737)(23.500,56.156)(24.667,55.608)(25.833,55.093)(27.000,54.608)(28.167,54.154)(29.333,53.729)(30.500,53.332)(31.667,52.963)(32.833,52.621)
(34.000,52.304)(35.167,52.012)(36.333,51.745)(37.500,51.500)(38.667,51.278)(39.833,51.077)(41.000,50.896)(42.167,50.735)(43.333,50.593)(44.500,50.468)(45.667,50.360)(46.833,50.269)
(48.000,50.192)(49.167,50.130)(50.333,50.081)(51.500,50.044)(52.667,50.019)(53.833,50.005)(55.000,50.000)(56.167,50.004)(57.333,50.017)(58.500,50.036)(59.667,50.062)(60.833,50.093)
(62.000,50.128)(63.167,50.167)(64.333,50.209)(65.500,50.252)(66.667,50.296)(67.833,50.341)(69.000,50.384)(70.167,50.426)(71.333,50.465)(72.500,50.500)(73.667,50.531)(74.833,50.557)
(76.000,50.576)(77.167,50.588)(78.333,50.593)(79.500,50.588)(80.667,50.574)(81.833,50.549)(83.000,50.512)(84.167,50.463)(85.333,50.401)(86.500,50.324)(87.667,50.232)(88.833,50.125)
(90.000,50.000)(91.167,49.858)(92.333,49.697)(93.500,49.516)(94.667,49.315)(95.833,49.093)(97.000,48.848)(98.167,48.580)(99.333,48.289)(100.500,47.972)(101.667,47.630)(102.833,47.261)
(104.000,46.864)(105.167,46.439)(106.333,45.985)(107.500,45.500)(108.667,44.984)(109.833,44.437)(111.000,43.856)(112.167,43.242)(113.333,42.593)(114.500,41.908)(115.667,41.187)(116.833,40.429)
(118.000,39.632)(119.167,38.796)(120.333,37.921)(121.500,37.004)(122.667,36.046)(123.833,35.045)(125.000,34.000)}

{\color{red}\linethickness{0.8pt}\polyline(20.000,53.200)(21.167,52.711)(22.333,52.257)(23.500,51.836)(24.667,51.448)(25.833,51.093)(27.000,50.768)(28.167,50.474)(29.333,50.209)(30.500,49.972)(31.667,49.763)(32.833,49.581)
(34.000,49.424)(35.167,49.292)(36.333,49.185)(37.500,49.100)(38.667,49.038)(39.833,48.997)(41.000,48.976)(42.167,48.975)(43.333,48.993)(44.500,49.028)(45.667,49.080)(46.833,49.149)
(48.000,49.232)(49.167,49.330)(50.333,49.441)(51.500,49.564)(52.667,49.699)(53.833,49.845)(55.000,50.000)(56.167,50.164)(57.333,50.337)(58.500,50.516)(59.667,50.702)(60.833,50.893)
(62.000,51.088)(63.167,51.287)(64.333,51.489)(65.500,51.692)(66.667,51.896)(67.833,52.101)(69.000,52.304)(70.167,52.506)(71.333,52.705)(72.500,52.900)(73.667,53.091)(74.833,53.277)
(76.000,53.456)(77.167,53.628)(78.333,53.793)(79.500,53.948)(80.667,54.094)(81.833,54.229)(83.000,54.352)(84.167,54.463)(85.333,54.561)(86.500,54.644)(87.667,54.712)(88.833,54.765)
(90.000,54.800)(91.167,54.818)(92.333,54.817)(93.500,54.796)(94.667,54.755)(95.833,54.693)(97.000,54.608)(98.167,54.500)(99.333,54.369)(100.500,54.212)(101.667,54.030)(102.833,53.821)
(104.000,53.584)(105.167,53.319)(106.333,53.025)(107.500,52.700)(108.667,52.344)(109.833,51.957)(111.000,51.536)(112.167,51.082)(113.333,50.593)(114.500,50.068)(115.667,49.507)(116.833,48.909)
(118.000,48.272)(119.167,47.596)(120.333,46.881)(121.500,46.124)(122.667,45.326)(123.833,44.485)(125.000,43.600)}

{\color{green!50!black}\linethickness{0.8pt}\polyline(20.000,48.400)(21.167,48.071)(22.333,47.777)(23.500,47.516)(24.667,47.288)(25.833,47.093)(27.000,46.928)(28.167,46.794)(29.333,46.689)(30.500,46.612)(31.667,46.563)(32.833,46.541)
(34.000,46.544)(35.167,46.572)(36.333,46.625)(37.500,46.700)(38.667,46.798)(39.833,46.917)(41.000,47.056)(42.167,47.215)(43.333,47.393)(44.500,47.588)(45.667,47.800)(46.833,48.029)
(48.000,48.272)(49.167,48.530)(50.333,48.801)(51.500,49.084)(52.667,49.379)(53.833,49.685)(55.000,50.000)(56.167,50.324)(57.333,50.657)(58.500,50.996)(59.667,51.342)(60.833,51.693)
(62.000,52.048)(63.167,52.407)(64.333,52.769)(65.500,53.132)(66.667,53.496)(67.833,53.861)(69.000,54.224)(70.167,54.586)(71.333,54.945)(72.500,55.300)(73.667,55.651)(74.833,55.997)
(76.000,56.336)(77.167,56.668)(78.333,56.993)(79.500,57.308)(80.667,57.614)(81.833,57.909)(83.000,58.192)(84.167,58.463)(85.333,58.721)(86.500,58.964)(87.667,59.192)(88.833,59.405)
(90.000,59.600)(91.167,59.778)(92.333,59.937)(93.500,60.076)(94.667,60.195)(95.833,60.293)(97.000,60.368)(98.167,60.420)(99.333,60.449)(100.500,60.452)(101.667,60.430)(102.833,60.381)
(104.000,60.304)(105.167,60.199)(106.333,60.065)(107.500,59.900)(108.667,59.704)(109.833,59.477)(111.000,59.216)(112.167,58.922)(113.333,58.593)(114.500,58.228)(115.667,57.827)(116.833,57.389)
(118.000,56.912)(119.167,56.396)(120.333,55.841)(121.500,55.244)(122.667,54.606)(123.833,53.925)(125.000,53.200)}

{\color{magenta}\linethickness{0.8pt}\polyline(20.000,62.800)(21.167,61.991)(22.333,61.217)(23.500,60.476)(24.667,59.768)(25.833,59.093)(27.000,58.448)(28.167,57.834)(29.333,57.249)(30.500,56.692)(31.667,56.163)(32.833,55.661)
(34.000,55.184)(35.167,54.732)(36.333,54.305)(37.500,53.900)(38.667,53.518)(39.833,53.157)(41.000,52.816)(42.167,52.495)(43.333,52.193)(44.500,51.908)(45.667,51.640)(46.833,51.389)
(48.000,51.152)(49.167,50.930)(50.333,50.721)(51.500,50.524)(52.667,50.339)(53.833,50.165)(55.000,50.000)(56.167,49.844)(57.333,49.697)(58.500,49.556)(59.667,49.422)(60.833,49.293)
(62.000,49.168)(63.167,49.047)(64.333,48.929)(65.500,48.812)(66.667,48.696)(67.833,48.581)(69.000,48.464)(70.167,48.346)(71.333,48.225)(72.500,48.100)(73.667,47.971)(74.833,47.837)
(76.000,47.696)(77.167,47.548)(78.333,47.393)(79.500,47.228)(80.667,47.054)(81.833,46.869)(83.000,46.672)(84.167,46.463)(85.333,46.241)(86.500,46.004)(87.667,45.752)(88.833,45.485)
(90.000,45.200)(91.167,44.898)(92.333,44.577)(93.500,44.236)(94.667,43.875)(95.833,43.493)(97.000,43.088)(98.167,42.660)(99.333,42.209)(100.500,41.732)(101.667,41.230)(102.833,40.701)
(104.000,40.144)(105.167,39.559)(106.333,38.945)(107.500,38.300)(108.667,37.624)(109.833,36.917)(111.000,36.176)(112.167,35.402)(113.333,34.593)(114.500,33.748)(115.667,32.867)(116.833,31.949)
(118.000,30.992)(119.167,29.996)(120.333,28.961)(121.500,27.884)(122.667,26.766)(123.833,25.605)(125.000,24.400)}

{\color{orange}\linethickness{0.8pt}\polyline(20.000,67.600)(21.167,66.631)(22.333,65.697)(23.500,64.796)(24.667,63.928)(25.833,63.093)(27.000,62.288)(28.167,61.514)(29.333,60.769)(30.500,60.052)(31.667,59.363)(32.833,58.701)
(34.000,58.064)(35.167,57.452)(36.333,56.865)(37.500,56.300)(38.667,55.758)(39.833,55.237)(41.000,54.736)(42.167,54.255)(43.333,53.793)(44.500,53.348)(45.667,52.920)(46.833,52.509)
(48.000,52.112)(49.167,51.730)(50.333,51.361)(51.500,51.004)(52.667,50.659)(53.833,50.325)(55.000,50.000)(56.167,49.684)(57.333,49.377)(58.500,49.076)(59.667,48.782)(60.833,48.493)
(62.000,48.208)(63.167,47.927)(64.333,47.649)(65.500,47.372)(66.667,47.096)(67.833,46.821)(69.000,46.544)(70.167,46.266)(71.333,45.985)(72.500,45.700)(73.667,45.411)(74.833,45.117)
(76.000,44.816)(77.167,44.508)(78.333,44.193)(79.500,43.868)(80.667,43.534)(81.833,43.189)(83.000,42.832)(84.167,42.463)(85.333,42.081)(86.500,41.684)(87.667,41.272)(88.833,40.845)
(90.000,40.400)(91.167,39.938)(92.333,39.457)(93.500,38.956)(94.667,38.435)(95.833,37.893)(97.000,37.328)(98.167,36.740)(99.333,36.129)(100.500,35.492)(101.667,34.830)(102.833,34.141)
(104.000,33.424)(105.167,32.679)(106.333,31.905)(107.500,31.100)(108.667,30.264)(109.833,29.397)(111.000,28.496)(112.167,27.562)(113.333,26.593)(114.500,25.588)(115.667,24.547)(116.833,23.469)
(118.000,22.352)(119.167,21.196)(120.333,20.001)(121.500,18.764)(122.667,17.486)(123.833,16.165)(125.000,14.800)}

\end{picture}
\par\small
\textcolor{blue}{$a=32$},\quad
\textcolor{red}{$a=31.9$},\quad
\textcolor{green!50!black}{$a=31.8$},\quad
\textcolor{magenta}{$a=32.1$},\quad
\textcolor{orange}{$a=32.2$}.
\end{center}
Every graph passes through $(1,0)$. At $a=32$, the root $1$ is
double, so its graph touches the axis there and also crosses at $2$.
For $a=31.9$, the three real crossings separate. For $a=31.8$,
two real roots move outside the plotted window; the graph within
$[0,3]$ shows only the crossing at $1$, even though all three roots
remain real. For $a=32.1,32.2$, the quadratic discriminant is negative,
so those two roots are genuinely nonreal, not merely off screen.
The coalescence threshold for that pair is $a=1537/48\approx32.020833$.
\answer{The exact factorization explains the changes. A graph on a
limited window alone cannot distinguish a root outside the window from
a nonreal root; the discriminant and computed values supply that distinction.}
